\section{Pencil realizations: hinges concurrent within a common panel}
\label{sec:pencil}

Katoh--Tanigawa's hinge-coplanar (panel) model
(\cref{def:coplanar-panel-realization}) assigns each body of a
body-hinge framework a hyperplane --- its \emph{panel} --- containing
every incident hinge. By projective duality
(\cref{sec:molecule-modelling-projective}) this coplanarity condition
is dual to hinge-\emph{concurrency}: every hinge at a body passing
through a common point. This section states the strengthening in
which both hold simultaneously at every body --- each body's hinges
lie in a common panel \emph{and} pass through a common point of that
panel, a pencil of lines through a point in a plane --- and proves
that the condition is itself projectively self-dual: transporting a
realization along the polarity of \cref{lem:panel-hinge-dual-molecular}
produces another realization of the same kind, with the roles of the
panel normal and the concurrency point exchanged. In the molecular
reading of \cref{sec:molecule-modelling-defs} (atoms as bodies, bond
lines as hinges, centres as concurrency points) the condition says
that every atom's incident bonds lie in a common plane through the
atom's centre. Whether every spanning multigraph admits such a
realization attaining the deficiency rank of
\cref{thm:theorem-55-6-multigraph} is open.

\subsection{Extensors through a point}
\label{sec:pencil-through-point}

The panel side of hinge-coplanarity is stated by containment: a screw
element lies in a hyperplane when it is the extensor of points of that
hyperplane (\cref{def:genuine-hinge-realization}). The dual notion ---
a screw element \emph{passing through} a point --- replaces
containment in the hyperplane by containment of the point in the
element's span.

\begin{definition}[Extensor through a point]
  \label{def:extensor-through-point}
  \lean{CombinatorialRigidity.Molecular.ExtensorThroughPoint}
  \leanok
  \uses{def:rigidity-matrix, def:extensor}
  A screw element $C \in \mathrm{ScrewSpace}\,k$ \emph{passes through}
  the homogeneous point $q \in K^{k+2}$ when $C$ is the extensor of $k$
  points $p_1, \dots, p_k \in K^{k+2}$ whose span contains $q$.
\end{definition}

This mirrors the witness shape of \cref{def:genuine-hinge-realization}
exactly, with span-containment of $q$ in place of orthogonality to a
normal: the degenerate $C = 0$ satisfies it (take every $p_i = q$),
so nonzero-ness of $C$ is carried as a separate conjunct wherever the
predicate is used below.

\subsection{Pencil realizations}
\label{sec:pencil-realization}

\begin{definition}[Pencil realization]
  \label{def:pencil-panel-realization}
  \lean{CombinatorialRigidity.Molecular.HasPencilPanelRealization}
  \leanok
  \uses{def:coplanar-panel-realization, def:extensor-through-point}
  A hinge-coplanar panel realization of a multigraph $G$
  (\cref{def:coplanar-panel-realization}), with normal assignment $v
  \mapsto n_v$, has a \emph{pencil realization} when it carries in
  addition a per-body homogeneous point $v \mapsto q_v \in K^{k+2}$
  such that every body of $V(G)$ has a nonzero point $q_v$, incident
  to its own panel ($q_v \cdot n_v = 0$), and every link's supporting
  extensor passes through the points of both its endpoints
  (\cref{def:extensor-through-point}).
\end{definition}

The incidence $q_v \cdot n_v = 0$ is forced at any body carrying a
genuine hinge --- the hinge's spanning points lie in $n_v^{\perp}$ and
contain $q_v$ in their span, so $q_v \in n_v^{\perp}$ --- and recording
it as a conjunct keeps an isolated body's data honest while making the
panel and point roles symmetric under the polarity below. Concurrency
of two coplanar lines is automatic, so at a body of degree at most two
the point conjunct adds nothing to hinge-coplanarity; the condition
first bites at a body whose hinges number three or more.

\subsection{The polarity and self-duality}
\label{sec:pencil-duality}

Fix $K = \R$ and grade $k = 2$ ($d = 3$), the setting of
\cref{lem:panel-hinge-dual-molecular}: the polarity there carries the
extensor of any two homogeneous points $v_0, v_1 \in \R^4$ to the
panel-support extensor $C(v_0,v_1)$ (\cref{def:panel-support-extensor}),
the join-of-two-points $\leftrightarrow$ meet-of-two-panels identity in
its full generality, not only for the homogenized poles of a molecular
framework. Two consequences carry the point-through and panel-in
conditions across the polarity.

\begin{lemma}[The point-join condition transports to the panel-meet side]
  \label{lem:pencil-transport-through-to-in}
  \lean{CombinatorialRigidity.Molecular.extensorInPanel_screwComplementIso_of_extensorThroughPoint}
  \leanok
  \uses{def:extensor-through-point, def:genuine-hinge-realization,
        lem:panel-hinge-dual-molecular}
  If a screw element $C \in \mathrm{ScrewSpace}\,2$ passes through the
  homogeneous point $q \in \R^4$ (\cref{def:extensor-through-point}),
  then its polarity image lies in the panel with normal $q$
  (\cref{def:genuine-hinge-realization}).
\end{lemma}
\begin{proof}
  \leanok
  \uses{lem:case-III-claim612-line-in-panel-union, def:panel-support-extensor}
  If $C = 0$ the image is $0$, which lies in every panel. Otherwise $C$
  is the extensor of an independent pair $v_0, v_1$ with $q$ in their
  span; the polarity (\cref{lem:panel-hinge-dual-molecular}) carries
  $C$ to the panel-support extensor $C(v_0,v_1)$
  (\cref{def:panel-support-extensor}), which by the point-join
  $\leftrightarrow$ panel-meet duality
  (\cref{lem:case-III-claim612-line-in-panel-union}) is the extensor of
  an independent pair $p_0', p_1'$ orthogonal to both $v_0$ and $v_1$.
  Since $q$ lies in the span of $v_0, v_1$ and the dot product is
  linear, $p_0'$ and $p_1'$ are also orthogonal to $q$, so the image
  lies in the panel with normal $q$.
\end{proof}

\begin{lemma}[The panel-meet condition transports to the point-join side]
  \label{lem:pencil-transport-in-to-through}
  \lean{CombinatorialRigidity.Molecular.extensorThroughPoint_screwComplementIso_of_extensorInPanel}
  \leanok
  \uses{def:extensor-through-point, def:genuine-hinge-realization,
        lem:panel-hinge-dual-molecular}
  If a screw element $C \in \mathrm{ScrewSpace}\,2$ lies in the panel
  with normal $q \in \R^4$ (\cref{def:genuine-hinge-realization}), then
  its polarity image passes through $q$
  (\cref{def:extensor-through-point}).
\end{lemma}
\begin{proof}
  \leanok
  \uses{lem:case-III-claim612-line-in-panel-union, def:panel-support-extensor}
  If $C = 0$ the image is $0$, which passes through every point (take
  the constant family at $q$). Otherwise $C$ is the extensor of an
  independent pair $v_0, v_1$ both orthogonal to $q$; the polarity
  carries $C$ to the panel-support extensor $C(v_0,v_1)$
  (\cref{def:panel-support-extensor}), which by
  \cref{lem:case-III-claim612-line-in-panel-union} is the extensor of
  an independent pair $p_0', p_1'$ orthogonal to both $v_0$ and $v_1$.
  The orthogonal complement of $\{v_0,v_1\}$ is $2$-dimensional, hence
  equal to the span of the independent pair $p_0', p_1'$; as $q$ also
  lies in this complement, $q \in \mathrm{span}\{p_0',p_1'\}$, i.e.\ the
  image passes through $q$.
\end{proof}

\begin{lemma}[Self-duality of the pencil condition]
  \label{lem:pencil-self-dual}
  \lean{CombinatorialRigidity.Molecular.hasPencilPanelRealization_mapExtensor_screwComplementIso}
  \leanok
  \uses{def:pencil-panel-realization, lem:panel-hinge-dual-molecular,
        thm:projective-invariance}
  Transporting a pencil realization $(F, n, q)$ of a multigraph $G$
  along the polarity of \cref{lem:panel-hinge-dual-molecular}
  (\cref{thm:projective-invariance}) produces a pencil realization
  $(\Lambda F, q, n)$ of $G$, with the roles of the normal and the
  point exchanged.
\end{lemma}
\begin{proof}
  \leanok
  \uses{lem:pencil-transport-through-to-in, lem:pencil-transport-in-to-through}
  The transported framework carries the same multigraph and nonzero
  supporting extensors, since the polarity is a linear automorphism
  (\cref{thm:projective-invariance}). Its new normal assignment is $q$,
  each nonzero by hypothesis, and the incidence $n_v \cdot q_v = 0$
  read the other way gives $q_v \cdot n_v = 0$, so it is preserved.
  Every link's transported supporting extensor lies in the panel of
  each endpoint's new normal $q$ by
  \cref{lem:pencil-transport-through-to-in} (the old point-join
  conditions become panel-meet conditions), and passes through each
  endpoint's new point $n$ by \cref{lem:pencil-transport-in-to-through}
  (the old panel-meet conditions become point-join conditions).
\end{proof}

\subsection{The coincident-panel pencil pair}
\label{sec:pencil-base}

Katoh--Tanigawa's two-vertex base realization of a parallel edge pair
(\cref{lem:theorem-55-base-producer-parallel}, Lemma~5.3, p.~670)
places both hinges in a single shared panel $n^{\perp}$ with two
linearly independent supporting extensors
(\cref{lem:extensor-pair-in-panel}), which fill the whole screw space
$K^6$ and make the two bodies rigid. On the pencil condition the same
construction carries a common concurrency point at no extra cost: two
lines of a flat pencil already share the point of the pencil.

\begin{lemma}[Two independent extensors through a common point of a panel]
  \label{lem:extensor-pair-through-point}
  \lean{CombinatorialRigidity.Molecular.exists_linearIndependent_extensor_pair_through_point}
  \leanok
  \uses{def:extensor-through-point, def:genuine-hinge-realization}
  For any normal $n \in K^4$ there is a nonzero point $q \in n^{\perp}$
  and two screw elements $C, C' \in \mathrm{ScrewSpace}\,2$, each lying
  in the panel with normal $n$ (\cref{def:genuine-hinge-realization})
  and passing through $q$ (\cref{def:extensor-through-point}), whose
  $2$-extensors are linearly independent.
\end{lemma}
\begin{proof}
  \leanok
  \uses{lem:extensor-pair-in-panel}
  The panel $n^{\perp}$ has dimension at least $3$ in $K^4$, so it
  carries three linearly independent vectors $q, a, b$. Take the two
  pencil lines $q \vee a$ and $q \vee b$ through $q$: both lie in
  $n^{\perp}$ and contain $q$ in their span, so each lies in the panel
  and passes through $q$. Their extensors are linearly independent
  because $\{q, a, b\}$ is --- the shared-vector case of the
  independence of extensors on two distinct point subsets behind
  \cref{lem:extensor-pair-in-panel}. These are the two distinct pencil
  lines through a common point of the shared panel whose stacked row
  spaces fill $K^6$, the rank mechanism of the two-body pencil base
  case.
\end{proof}

The base region of Katoh--Tanigawa's induction is realized directly at
one and two bodies (\cref{lem:theorem-55-base-producer}); the pencil
condition first bites at a body of degree three or more, so on two
bodies the only content is the parallel-pair case, where the pencil pin
is met for free by the pair above.

\begin{lemma}[Two-body coincident-panel pencil realization]
  \label{lem:pencil-base-parallel-pair}
  \lean{CombinatorialRigidity.Molecular.exists_pencilPanelRealization_parallel_pair}
  \leanok
  \uses{def:pencil-panel-realization, def:rank-hypothesis}
  Let $G$ be a two-vertex minimal-$0$-dof multigraph realizing the
  parallel-pair case (\cref{lem:theorem-55-base-producer-parallel}):
  $V(G) = \{x, y\}$ with $x \ne y$ and $E(G) = \{e, f\}$ two distinct
  edges both joining $x$ and $y$. Then $G$ has a pencil realization
  (\cref{def:pencil-panel-realization}) that is infinitesimally rigid on
  its two bodies $V(G) = \{x, y\}$ (\cref{def:rank-hypothesis}, the
  $V(G)$-relative rank $D$ of the $\mathrm{def}(\tilde G) = 0$ case).
\end{lemma}
\begin{proof}
  \leanok
  \uses{lem:extensor-pair-through-point, lem:theorem-55-base}
  Place both bodies on a single panel $n^{\perp}$ at a fixed nonzero
  normal $n$ with a single shared concurrency point $q \in n^{\perp}$,
  and assign the two parallel hinges the two independent pencil lines
  through $q$ of \cref{lem:extensor-pair-through-point}; every edge label
  off $\{e, f\}$ reuses the first hinge, so all supporting extensors are
  nonzero. Each hinge lies in the common panel and passes through the
  common point, so the data is a pencil realization. The two independent
  supporting extensors give the combined hinge-row blocks full rank on
  the relative screw $S(x) - S(y)$, so the framework is infinitesimally
  rigid on $\{x, y\} = V(G)$ (\cref{lem:theorem-55-base}).
\end{proof}

\subsection{Pencil cycles: degree-two concurrency is automatic}
\label{sec:pencil-cycle}

Katoh--Tanigawa's cycle realization (\cref{lem:cycle-realization},
Lemma~5.4) places a cycle of bodies with each hinge the meet of two
adjacent panels. Every body has degree two, and at a body of degree at
most two the pencil concurrency pin costs nothing: two lines lying in a
common plane always meet in a point of that plane. The following records
this projective fact in the containment model, so that the cycle
realization --- indeed any hinge-coplanar realization
(\cref{def:coplanar-panel-realization}) restricted to its degree-two
bodies --- already carries the per-body concurrency point a pencil
realization asks for.

\begin{lemma}[Two coplanar hinges share a concurrency point]
  \label{lem:coplanar-hinges-concurrent}
  \lean{CombinatorialRigidity.Molecular.exists_concurrency_point_of_extensorInPanel_pair}
  \leanok
  \uses{def:genuine-hinge-realization, def:extensor-through-point}
  Let $n \in K^4$ be nonzero and let $C_1, C_2 \in \mathrm{ScrewSpace}\,2$
  be nonzero screw elements each lying in the panel with normal $n$
  (\cref{def:genuine-hinge-realization}). Then there is a nonzero point
  $q \in n^{\perp}$ through which both pass
  (\cref{def:extensor-through-point}).
\end{lemma}
\begin{proof}
  \leanok
  Each hinge is the extensor of an independent pair, spanning a
  $2$-dimensional subspace of the panel $n^{\perp}$, which has dimension
  $3$ in $K^4$. Two $2$-dimensional subspaces of a $3$-dimensional space
  meet in dimension at least $2 + 2 - 3 = 1$ (the modular law), so their
  intersection contains a nonzero point $q$. Lying in each hinge's span,
  $q$ passes through both hinges; and lying in $n^{\perp}$ it is
  orthogonal to $n$.
\end{proof}

The cycle itself carries a hinge-coplanar realization, rigid on all of
its bodies, built directly from the cyclic family of shared panel normals
rather than through the meet model of the original cycle realization.

\begin{lemma}[Coplanar realization of a cycle]
  \label{lem:cycle-coplanar-realization}
  \lean{CombinatorialRigidity.Molecular.exists_coplanarPanelRealization_cycle}
  \leanok
  \uses{def:cycle-data, def:coplanar-panel-realization, def:rank-hypothesis}
  Let $G$ be presented as a cycle of $m$ bodies (\cref{def:cycle-data})
  with $3 \le m \le 4$. Then $G$ has a hinge-coplanar panel realization
  (\cref{def:coplanar-panel-realization}) that is infinitesimally rigid on
  all of its bodies $V(G)$ (\cref{def:rank-hypothesis}).
\end{lemma}
\begin{proof}
  \leanok
  \uses{lem:cycle-normals, lem:cycle-realization-rigid,
        def:panel-support-extensor, def:genuine-hinge-realization}
  At grade $k = 2$ the cyclic shared-normal family
  (\cref{lem:cycle-normals}) supplies $m$ panel normals $n_0, \dots,
  n_{m-1}$ --- the range $3 \le m \le 4 = k + 2$ is exactly its
  hypothesis --- with each consecutive pair $n_i, n_{i+1}$ linearly
  independent. The $i$-th cycle edge is assigned the panel meet
  $C(n_i, n_{i+1})$ (\cref{def:panel-support-extensor}), nonzero and lying
  in both endpoint panels $n_i^{\perp}$, $n_{i+1}^{\perp}$
  (\cref{def:genuine-hinge-realization}); every edge label off the cycle
  is given a fixed nonzero fallback, so all supporting extensors are
  nonzero. Body $i$ carries normal $n_i$, and both hinges incident to it
  lie in $n_i^{\perp}$, so the data is a hinge-coplanar realization. The
  $m$ cyclic meets are linearly independent (\cref{lem:cycle-normals}), so
  the framework is infinitesimally rigid on the cycle's bodies, which are
  all of $V(G)$ (\cref{lem:cycle-realization-rigid}).
\end{proof}

The cycle length range here is honest rather than triangle-only: a cycle
is at least a triangle ($m \ge 3$), and the shared-normal family exists
only up to $m \le k + 2 = 4$ at $d = 3$, so both the triangle and the
quadrilateral are realized. Since every body of a cycle has degree two,
the concurrency fact above (\cref{lem:coplanar-hinges-concurrent})
supplies each body's point, upgrading the coplanar realization to a full
pencil realization.

\begin{lemma}[Pencil realization of a cycle]
  \label{lem:cycle-pencil-realization}
  \lean{CombinatorialRigidity.Molecular.exists_pencilPanelRealization_cycle}
  \leanok
  \uses{def:cycle-data, def:pencil-panel-realization, def:rank-hypothesis}
  Let $G$ be presented as a cycle of $m$ bodies (\cref{def:cycle-data})
  with $3 \le m \le 4$. Then $G$ has a pencil realization
  (\cref{def:pencil-panel-realization}) that is infinitesimally rigid on
  all of its bodies $V(G)$ (\cref{def:rank-hypothesis}).
\end{lemma}
\begin{proof}
  \leanok
  \uses{lem:cycle-coplanar-realization, lem:coplanar-hinges-concurrent,
        def:panel-support-extensor, def:extensor-through-point}
  Take the coplanar realization of the cycle
  (\cref{lem:cycle-coplanar-realization}), whose $i$-th edge is the panel
  meet $C(n_i, n_{i+1})$ (\cref{def:panel-support-extensor}). Body $i$ has
  exactly two incident hinges, $C(n_{i-1}, n_i)$ and $C(n_i, n_{i+1})$,
  both lying in the panel $n_i^{\perp}$; by
  \cref{lem:coplanar-hinges-concurrent} they share a nonzero point
  $q_i \in n_i^{\perp}$, through which both pass
  (\cref{def:extensor-through-point}). Assign body $i$ the point $q_i$.
  Then each point is nonzero and incident to its own panel, and the
  edge $C(n_i, n_{i+1})$ passes through the points of both its endpoints:
  through $q_i$ as body $i$'s second hinge, and through $q_{i+1}$ as body
  $i+1$'s first hinge (its first incident hinge is $C(n_i, n_{i+1})$
  since $(i+1) - 1 = i$). Rigidity is inherited from the coplanar
  realization.
\end{proof}

\subsection{The two-pencil extension step}
\label{sec:pencil-extension}

Katoh--Tanigawa's induction extends a hinge-coplanar realization across
a re-added edge by supplying a supporting extensor in the meet of the
two endpoint panels (p.~670): for any two normals the common perp
$n_u^{\perp} \cap n_v^{\perp} \subseteq K^4$ has dimension at least
$4 - 2 = 2$, so it carries a nonzero hinge lying in both panels, with
no transversality of the normals needed. On the pencil condition the
re-added hinge must additionally pass through the concurrency point of
each endpoint. Where in the coplanar model a common hinge always
exists, the pencil model asks for one through two prescribed points as
well, and this constrains the two bodies: a hinge lying in both panels
and passing through both points confines each point to the common perp,
so each concurrency point must lie not only in its own panel but in the
other body's panel too. These two cross-incidences are exactly what the
following existence step assumes --- making explicit the obligation the
coplanar strip-and-re-add move leaves implicit.

\begin{lemma}[Two-pencil extension hinge]
  \label{lem:two-pencil-extension}
  \lean{CombinatorialRigidity.Molecular.exists_extensor_two_pencils}
  \leanok
  \uses{def:genuine-hinge-realization, def:extensor-through-point}
  Let $n_u, n_v \in K^4$ be panel normals and $q_u, q_v \in K^4$
  concurrency points with $q_u$ incident to its own panel
  ($q_u \cdot n_u = 0$) and $q_v$ to its own ($q_v \cdot n_v = 0$),
  and suppose $q_u \ne 0$. If the two cross-incidences
  $q_u \cdot n_v = 0$ and $q_v \cdot n_u = 0$ hold --- each
  concurrency point lies in the other body's panel --- then there is a
  nonzero $C \in \mathrm{ScrewSpace}\,2$ lying in both panels
  (\cref{def:genuine-hinge-realization}) and passing through both points
  (\cref{def:extensor-through-point}).
\end{lemma}
\begin{proof}
  \leanok
  Under the four incidences both points lie in the common perp
  $W = n_u^{\perp} \cap n_v^{\perp}$, which has dimension at least $2$
  in $K^4$ regardless of whether the normals are independent. It
  suffices to exhibit a $2$-dimensional subspace of $W$ containing both
  points; its extensor is then a nonzero hinge in both panels through
  both points. If $q_u$ and $q_v$ are linearly independent their span
  is such a subspace. Otherwise $q_v$ lies on the line through $q_u$,
  and since $W$ has dimension at least $2$ while that line has dimension
  $1$, $W$ contains a vector $w$ independent from $q_u$; the span of
  $q_u$ and $w$ is a $2$-dimensional subspace of $W$ containing $q_u$,
  hence also $q_v$.
\end{proof}

The two cross-incidences are not only sufficient but necessary, so the
extension hinge exists precisely when the two bodies' concurrency points
are mutually panel-incident.

\begin{lemma}[Two-pencil extension biconditional]
  \label{lem:two-pencil-extension-iff}
  \lean{CombinatorialRigidity.Molecular.exists_extensor_two_pencils_iff}
  \leanok
  \uses{lem:two-pencil-extension, def:genuine-hinge-realization,
        def:extensor-through-point}
  With the hypotheses of \cref{lem:two-pencil-extension} --- the points
  $q_u, q_v$ incident to their own panels and $q_u \ne 0$ --- the two
  cross-incidences $q_u \cdot n_v = 0$ and $q_v \cdot n_u = 0$ hold if and
  only if there is a nonzero screw element lying in both panels
  (\cref{def:genuine-hinge-realization}) and passing through both points
  (\cref{def:extensor-through-point}).
\end{lemma}
\begin{proof}
  \leanok
  \uses{lem:two-pencil-extension}
  The forward implication is \cref{lem:two-pencil-extension}. For the
  converse, suppose such a nonzero hinge $C$ exists. A nonzero decomposable
  $2$-extensor determines its plane: two point pairs with the same nonzero
  $2$-extensor span the same plane (Plücker injectivity, the converse of the
  point-pair proportionality of \cref{lem:case-III-claim612-line-in-panel-union}).
  So the points through which $C$ passes and the points spanning it inside a
  panel span one and the same plane; since $C$ passes through $q_u$ and lies
  in the panel of $n_v$, the point $q_u$ lies in a span of vectors orthogonal
  to $n_v$, giving $q_u \cdot n_v = 0$, and symmetrically $q_v \cdot n_u = 0$.
\end{proof}

\subsection{Reduction on all spanning multigraphs}
\label{sec:pencil-reduction}

Katoh--Tanigawa's reduction of $0$-dof multigraphs (KT Theorem 4.9, restated
in \cref{sec:molecular-induction} without reference to the pencil condition)
proceeds by strong induction on $|V(G)|$ and at each step uses that the
multigraph in question is \emph{minimal}: no edge can be deleted without
changing the deficiency. Minimality is not available once an edge already
deleted must be reinserted with its supporting extensor constrained to pass
through two prescribed points (\cref{lem:two-pencil-extension}): the
reinserted edge can then fail to contribute any rank, once its endpoints'
concurrency points force every hinge extensor of the deleted graph into a
single three-dimensional subspace. The reduction is therefore restated here
on every spanning multigraph directly, with one additional combinatorial
fact removing the need for minimality: a loopless multigraph with every
vertex of degree at least three already has a proper rigid subgraph
(\cref{lem:pencil-min-degree-rigid}), so a graph without one has a vertex of
degree at most two, and two-edge-connectivity forces that degree to be
exactly two.

\begin{definition}[Pencil realization at the deficiency rank]
  \label{def:pencil-rank-hypothesis}
  \lean{CombinatorialRigidity.Molecular.HasPencilRealization}
  \leanok
  \uses{def:pencil-panel-realization, def:D-deficiency}
  A \emph{pencil realization at the deficiency rank} of a multigraph $G$ is
  a pencil realization $(F, n, q)$ of $G$
  (\cref{def:pencil-panel-realization}) whose rigidity rows span dimension
  \[
    \dim \operatorname{span}(\text{rigidity rows of } F)
      = D(|V(G)| - 1) - \operatorname{def}(\tilde G),
  \]
  the target of \cref{thm:theorem-55-6-multigraph} restricted to the pencil
  condition.
\end{definition}

\begin{lemma}[Proper rigid subgraph from a minimum degree of three]
  \label{lem:pencil-min-degree-rigid}
  \lean{Graph.exists_isProperRigidSubgraph_of_three_le_degree}
  \leanok
  \uses{def:rigid-subgraph, def:matroid-MG}
  Let $G$ be a loopless multigraph with $|V(G)| \ge 3$ and $D = \binom{d+1}{2}
  \ge 4$. If every vertex of $G$ has degree at least three, $G$ has a proper
  rigid subgraph.
\end{lemma}
\begin{proof}
  \leanok
  Let $v$ be a vertex of minimum degree $\delta \ge 3$. The handshake
  identity gives $\delta\,|V(G)| \le 2\,|E(G)|$, so the edges avoiding $v$
  number $|E(G)| - \delta \ge \delta(|V(G)|-2)/2$, and the $(D-1)$-fold
  fiber of these edges in $M(\tilde G)$ has size $(D-1)(|E(G)| - \delta) >
  D(|V(G)| - 2)$ once $D \ge 4$ --- strictly more than the $(D,D)$-sparsity
  bound on any vertex set omitting $v$. The fiber is therefore dependent in
  $M(\tilde G)$ and contains a circuit, whose induced subgraph is rigid
  (\cref{lem:circuit-induces-rigid}), spans at least two vertices, and
  omits $v$, hence is proper.
\end{proof}

\begin{theorem}[Reduction on all spanning multigraphs; cf.\ KT Theorem 4.9]
  \label{thm:pencil-reduction}
  \lean{Graph.pencil_reduction}
  \leanok
  \uses{def:cut-edges-2ec, def:rigid-subgraph, lem:pencil-min-degree-rigid}
  A property $P$ of multigraphs that is closed under five cases holds of
  every nonempty multigraph. The five cases are: \textup{(i)} $G$ has a
  loop (delete it and reinsert it); \textup{(ii)} $G$ is loopless with
  $|V(G)| \le 2$ (base case); \textup{(iii)} $G$ is loopless, $|V(G)| \ge
  3$, and not two-edge-connected (cut-edge decomposition,
  \cref{def:cut-edges-2ec}); \textup{(iv)} $G$ is loopless, $|V(G)| \ge 3$,
  and has a proper rigid subgraph (\cref{def:rigid-subgraph}); \textup{(v)}
  $G$ is loopless, two-edge-connected, $|V(G)| \ge 3$, has no proper rigid
  subgraph, and (by \cref{lem:pencil-min-degree-rigid}) some vertex has
  degree exactly two. Cases (i), (iii), (iv), (v) are handed the property
  at every strictly smaller graph in the lexicographic order on
  $(|V(G)|, |E(G)|)$.
\end{theorem}
\begin{proof}
  \leanok
  Strong induction on $(|V(G)|, |E(G)|)$, lexicographically: deleting a
  loop strictly decreases $|E(G)|$ at fixed $|V(G)|$, and every other case
  strictly decreases $|V(G)|$. Case (v)'s degree-two vertex follows from
  \cref{lem:pencil-min-degree-rigid} (no proper rigid subgraph forces a
  vertex of degree at most two) and two-edge-connectivity (degree at least
  two).
\end{proof}

\begin{lemma}[Simplicity from looplessness and no proper rigid subgraph]
  \label{lem:pencil-simple-of-noRigid}
  \lean{Graph.simple_of_loopless_of_noRigid}
  \leanok
  \uses{def:rigid-subgraph}
  A parallel edge pair joining two vertices $x, y$ of a multigraph $G$ with
  $|V(G)| \ge 3$ already forms a proper rigid subgraph on $\{x,y\}$, so
  case~(v) of \cref{thm:pencil-reduction} may assume $G$ simple.
\end{lemma}
\begin{proof}
  \leanok
  The two-vertex induced subgraph on a parallel pair is trivially rigid
  (KT Lemma 6.2) and its vertex set is a proper subset once $|V(G)| \ge 3$,
  so the absence of a proper rigid subgraph forces $G$ to have no parallel
  pair; combined with looplessness this is simplicity.
\end{proof}

\begin{lemma}[The loop case]
  \label{lem:pencil-loop-case}
  \lean{CombinatorialRigidity.Molecular.hasPencilRealization_of_isLoopAt}
  \leanok
  \uses{def:pencil-rank-hypothesis, lem:two-pencil-extension}
  Let $e$ be a loop of $G$ at a vertex $v$, and suppose $G - e$ has a
  pencil realization at the deficiency rank. Then so does $G$: assign $e$
  any pencil line of $v$'s own panel through $v$'s own point
  (\cref{lem:two-pencil-extension} at $n_u = n_v$, $q_u = q_v$), which
  contributes no rank, since $\operatorname{def}(\tilde G) =
  \operatorname{def}(\widetilde{G - e})$.
\end{lemma}
\begin{proof}
  \leanok
  Immediate from \cref{lem:two-pencil-extension} applied to $v$ against
  itself and the deficiency identity for deleting a loop.
\end{proof}

\begin{lemma}[Projective transport of a pencil realization]
  \label{lem:pencil-projective-transport}
  \lean{CombinatorialRigidity.Molecular.hasPencilPanelRealization_mapSupport_screwEquivOfLinearEquiv,
        CombinatorialRigidity.Molecular.extensorInPanel_screwEquivOfLinearEquiv,
        CombinatorialRigidity.Molecular.extensorThroughPoint_screwEquivOfLinearEquiv}
  \leanok
  \uses{def:pencil-panel-realization}
  Let $g$ be a linear automorphism of $K^4$ and $h$ a contragredient of it,
  meaning $\langle g x, h y\rangle = \langle x, y\rangle$ for the standard
  pairing. Transporting a pencil panel realization $(F, n, q)$ by acting on
  hinges through the induced automorphism of the screw space, on the
  concurrency points $q$ by $g$, and on the panel normals $n$ by $h$, again
  gives a pencil panel realization on the same multigraph. Its rigidity-row
  span has the same dimension, so a pencil realization at the deficiency
  rank transports to one.
\end{lemma}
\begin{proof}
  \leanok
  Panel containment and through-point incidence transport because the
  induced screw automorphism sends the extensor of spanning points
  $p_0, \dots$ to the extensor of $g p_0, \dots$; the contragredient
  identity turns each incidence $\langle p_i, n\rangle = 0$ into $\langle g
  p_i, h n\rangle = 0$, and $g$ carries the point-span containing $q$ to
  the one containing $g q$. The panel--point incidence $\langle q,
  n\rangle = 0$ is likewise preserved, and the automorphisms send nonzero
  data to nonzero data. Rank invariance is the change-of-screw-coordinates
  identity (this is the general-field companion of
  \cref{lem:pencil-self-dual}, whose polarity is one such automorphism).
\end{proof}

\begin{lemma}[Existence of a repositioning automorphism]
  \label{lem:pencil-cut-nondegeneracy}
  \lean{CombinatorialRigidity.Molecular.exists_reposition_cross_incidences,
        CombinatorialRigidity.Molecular.exists_contragredient_linearEquiv,
        CombinatorialRigidity.Molecular.exists_perp_linearIndependent}
  \leanok
  \uses{lem:two-pencil-extension-iff}
  Fix the panel/point data of the two sides at a crossing edge $uv$: normal
  $n_1(u)$ and point $pt_1(u) \ne 0$ on the $V_1$ side, normal $n_2(v)$ and
  point $pt_2(v) \ne 0$ on the $V_2$ side. Then there is a linear
  automorphism $g$ of $K^4$ with a contragredient $h$ ($\langle g x, h
  y\rangle = \langle x, y\rangle$ for the standard pairing) meeting the two
  cross-incidences of \cref{lem:two-pencil-extension-iff} once the $V_2$
  side is transported by $(g, h)$: $\langle pt_1(u), h\, n_2(v)\rangle = 0$
  and $\langle g\, pt_2(v), n_1(u)\rangle = 0$.
\end{lemma}
\begin{proof}
  \leanok
  Explicit, over any field. Every automorphism of $K^4$ has a
  contragredient, the inverse-transpose. Pick $a \in n_1(u)^\perp$ with
  $\{a, pt_1(u)\}$ independent and $b \in n_2(v)^\perp$ with $\{b,
  pt_2(v)\}$ independent, possible because each hyperplane $\perp$ is at
  least $3$-dimensional and so is not contained in a line. Let $g$ be the
  frame automorphism sending $pt_2(v) \mapsto a$ and $b \mapsto pt_1(u)$.
  Then $g\, pt_2(v) = a \in n_1(u)^\perp$ is the second incidence, and $g\,
  b = pt_1(u)$ with the contragredient identity turns the first into $b \in
  n_2(v)^\perp$.
\end{proof}

\begin{lemma}[The cut-edge case]
  \label{lem:pencil-cut-case}
  \lean{CombinatorialRigidity.Molecular.hasPencilRealization_of_not_twoEdgeConnected}
  \leanok
  \uses{def:pencil-rank-hypothesis, def:cut-edges-2ec,
        lem:two-pencil-extension-iff, lem:pencil-projective-transport,
        lem:pencil-cut-nondegeneracy}
  Let $G$ be not two-edge-connected, with at
  most one edge crossing some nonempty proper vertex set $V_1 \subsetneq
  V(G)$, $V_2 = V(G) \setminus V_1$. If $G[V_1]$ and $G[V_2]$ each have a
  pencil realization at the deficiency rank, so does $G$: place each side's
  realization independently and, if a crossing edge remains, reposition one
  side's frame along a projective automorphism so its two prescribed points
  meet the crossing edge's two cross-incidence conditions
  (\cref{lem:two-pencil-extension-iff}).
\end{lemma}
\begin{proof}
  \leanok
  \uses{lem:pencil-projective-transport, lem:pencil-cut-nondegeneracy,
        lem:cut-edge-decomposition, lem:two-pencil-extension}
  As in the cut-edge case of the panel-only induction
  (\cref{lem:case-cut-edge-realization}), with the crossing edge's
  supporting extensor supplied by \cref{lem:two-pencil-extension} once the
  two cross-incidences hold. The projective repositioning of one side is
  \cref{lem:pencil-projective-transport} (which preserves the deficiency
  rank); the automorphism meeting the two cross-incidences is furnished by
  \cref{lem:pencil-cut-nondegeneracy}.
\end{proof}

\begin{lemma}[The base case, $|V(G)| \le 2$]
  \label{lem:pencil-base-case}
  \uses{def:pencil-rank-hypothesis, lem:pencil-base-parallel-pair}
  Every loopless multigraph with at most two vertices has a pencil
  realization at the deficiency rank: the one-vertex, no-edge, and
  single-edge cases are immediate, and the parallel-edge-class case is
  \cref{lem:pencil-base-parallel-pair}.
\end{lemma}
\begin{proof}
  The one-vertex and edgeless cases contribute no rank; a single edge or a
  class of parallel edges between the two vertices reduces to
  \cref{lem:pencil-base-parallel-pair}, whose realization is
  infinitesimally rigid on both bodies.
\end{proof}

\begin{lemma}[Deficiency is unchanged by contracting a proper rigid subgraph]
  \label{lem:pencil-contraction-deficiency}
  \lean{Graph.rigidContract_deficiency_eq}
  \leanok
  \uses{def:rigid-subgraph, lem:contract-deficiency-conservation}
  Let $H$ be a proper rigid subgraph of $G$. Contracting $H$ to a single
  vertex leaves the deficiency unchanged,
  $\operatorname{def}(\widetilde{G / E(H)}) = \operatorname{def}(\tilde
  G)$, without assuming $G$ minimal.
\end{lemma}
\begin{proof}
  \leanok
  The matroid-contraction rank identity and the rank/deficiency bridge
  behind \cref{lem:contract-deficiency-conservation} already use no
  minimality; only the base/fiber-meeting half of the corresponding
  minimality-tracking statement did, and that half is not needed here.
\end{proof}

\begin{theorem}[Pencil realization from the contraction and split cases]
  \label{thm:pencil-conditional-realization}
  \uses{thm:pencil-reduction, def:pencil-rank-hypothesis, def:rigid-subgraph}
  Suppose that for every loopless multigraph $G$ with $|V(G)| \ge 3$, if
  every strictly smaller such graph has a pencil realization at the
  deficiency rank, then so does $G$ whenever it has a proper rigid
  subgraph, and so does $G$ whenever it is two-edge-connected with no
  proper rigid subgraph. Then every multigraph on the whole ambient body
  set has a pencil realization at the deficiency rank.
\end{theorem}
\begin{proof}
  \uses{thm:pencil-reduction, lem:pencil-loop-case, lem:pencil-base-case,
        lem:pencil-cut-case}
  Apply \cref{thm:pencil-reduction}: the loop case is
  \cref{lem:pencil-loop-case}, the base case is
  \cref{lem:pencil-base-case}, and the cut-edge case is
  \cref{lem:pencil-cut-case}; the contraction and split cases are the two
  hypotheses.
\end{proof}

\begin{fmlnote}
  \label{fmlnote:pencil-conditional-bare}
  The two hypotheses above are stated for the bare existence predicate
  (\cref{def:pencil-rank-hypothesis}). Discharging them is expected to need
  the induction hypothesis strengthened by a pencil-generic conjunct, the
  way Theorem~5.5's induction hypothesis carries a generic full-rank
  conjunct alongside the bare panel realization (\cref{def:rank-hypothesis}).
  Pinning that conjunct is the first task of the genericity argument these
  two cases are expected to need.
\end{fmlnote}
